\section{Analysis}

\startframedtask
On page 46

Analyze the imagery used in Oberon's description of Titania's resting place and his plan of action.

Point out the effect which is created (ll. 248 - 258)
\stopframedtask


\subsection{Analysis of Imagery}

\startitemize[packed]

\item {\bf Natural Imagery}: Oberon describes Titania’s resting place using rich and delicate natural imagery. Words like “wild thyme,” “oxlips,” “violet,” and “musk-roses” create a harmless and beautiful surrounding.  
\item {\bf Sensoring}: To describe the resting place Oberon not only uses Visual markers, but also other sensorign Impressions, like smell "sweet musk-roses." or taste.

\stopitemize

In Act II, Scene 1 of A Midsummer Night's Dream, lines 248-258, Oberon delivers
a speech that serves a dual purpose: it both describes Titania’s resting place
and outlines his intention to enchant her. In this speech Shakespeare heavily
relies on using imagery, contrast to deepen the audience’s understanding of
characters and themes in the play. Oberon’s language  and imagery is also used
to setting up a twist that shifts the tone from peaceful to manipulative and
dark.

The speech begins with a detailed description of Titania's resting place in the
woods. This speech evokes an image of a idyllic, safe, and almost magical
location.

In order to achieve that, Oberon uses a lot of nature connected words and
points of all the flowers surrounding the place like "violet", "woodbine",
"musk-roses" or "eglantine". Flowers represent the joy of life and the beauty
of nature. Given that they surround the place amplified the bond with nature. 

Also, Oberon creates an image of safety and comfort while talking about
Titianas resting place. For an example that Titiana is lulled in flowers (l.
254) within her over-canopied bank (l. 249-251). This supports the twist
between the evil and the comforting place in his speech.

In contrast to the beautiful nature mentioned, Oberon still has an evil plan
going on. For an instance he wants to use the magical flower to give Titania
hateful and torturing imaginations (l. 258). This causes a contrast between how
Oberon illustrates Titanias lovely and cute living space which makes her feel
safe and secure while still continueing his evil plan to make her fall in love
with a forest animal.

\blank
\hrule
\blank

{\bf The Twist:} starts with the snake\sidenote{Snakes are venemous, so evil}

\startitemize[m,packed]
\item A snake that is throwing of her skin, is changeging. This also stands for the twist and the change of tone in the text
\item The snake stands for oberon, planing behind titanias back
\stopitemize
